\documentclass{article}
\usepackage{amsmath}
\usepackage{amssymb}
\usepackage{amsfonts}
\usepackage{bm} % For bold math symbols like vectors

\newcommand{\mat}[1]{\mathbf{#1}} % Matrix
\newcommand{\vecbf}[1]{\bm{#1}} % Vector (bold)


\begin{document}

\section*{1. High-Correlation Feature Filter}
\subsection*{Initialization}
\begin{itemize}
    \item Correlation magnitude upper limit $\tau_{corr} \in [0, 1]$.
\end{itemize}

\subsection*{Fitting Process}
\paragraph{Inputs:} Data matrix $\mat{X} \in \mathbb{R}^{N \times D}$.
\paragraph{Procedure:}
\begin{enumerate}
    \item Compute $D \times D$ Pearson correlation matrix $\mat{C}$:
    $$ C_{ij} = \frac{\sum_{k=1}^{N} (X_{ki} - \bar{X}_i)(X_{kj} - \bar{X}_j)}{\sqrt{\sum_{k=1}^{N} (X_{ki} - \bar{X}_i)^2} \sqrt{\sum_{k=1}^{N} (X_{kj} - \bar{X}_j)^2}} $$
    \item Identify set of feature indices for removal $\mathcal{I}_{drop}$:
    $$ \mathcal{I}_{drop} = \{ j \mid \exists i < j \text{ such that } |C_{ij}| > \tau_{corr} \} $$
\end{enumerate}

\subsection*{Transformation Process}
\paragraph{Inputs:} Data matrix $\mat{X}' \in \mathbb{R}^{N' \times D'}$, set $\mathcal{I}_{drop}$.
\paragraph{Procedure:} Remove columns from $\mat{X}'$ with indices in $\mathcal{I}_{drop}$ (if valid).

\section*{2. Sequential Preprocessing Transformer}
A pipeline applying transformations sequentially: $\mat{X} \xrightarrow{T_1} \mat{X}^{(1)} \xrightarrow{T_2} \mat{X}^{(2)} \xrightarrow{T_3} \mat{X}^{(3)} \xrightarrow{T_4} \mat{X}^{(4)}$.

\subsection*{Stage 1: Low-Variance Feature Removal ($T_1$)}
\begin{itemize}
    \item For each feature $\vecbf{x}_j$ in input $\mat{X}^{(0)}$, calculate sample variance $s_j^2 = \frac{1}{N-1}\sum_{k=1}^{N} (X_{kj}^{(0)} - \bar{X}_j^{(0)})^2$.
    \item Remove feature $\vecbf{x}_j$ if $s_j^2 < \tau_{var}$ (e.g., $\tau_{var} = 0.05$).
    \item Output: $\mat{X}^{(1)}$.
\end{itemize}

\subsection*{Stage 2: High-Correlation Feature Filtering ($T_2$)}
\begin{itemize}
    \item Applies the filter from Section 1 to $\mat{X}^{(1)}$ (e.g., with $\tau_{corr} = 0.95$).
    \item Output: $\mat{X}^{(2)}$.
\end{itemize}

\subsection*{Stage 3: Data Scaling (Standardization) ($T_3$)}
\begin{itemize}
    \item For each feature $\vecbf{x}_j^{(2)}$ in $\mat{X}^{(2)}$ with mean $\bar{X}_j^{(2)}$ and std. dev. $s_j^{(2)}$:
    $$ X_{kj}^{(3)} = \frac{X_{kj}^{(2)} - \bar{X}_j^{(2)}}{s_j^{(2)}} $$
    (If $s_j^{(2)}=0$, $X_{kj}^{(3)}=0$).
    \item Output: $\mat{X}^{(3)}$.
\end{itemize}

\subsection*{Stage 4: Principal Component Analysis (PCA) ($T_4$)}
\begin{itemize}
    \item Input $\mat{X}^{(3)}$ (assumed centered or centered internally).
    \item Compute covariance matrix $\mat{S}_X = \frac{1}{N-1} (\mat{X}^{(3)})^T \mat{X}^{(3)}$.
    \item Perform eigen-decomposition of $\mat{S}_X$ to get eigenvalues $\lambda_i$ and eigenvectors $\vecbf{w}_i$.
    \item Select $D_{pca}$ eigenvectors $\vecbf{w}_1, \dots, \vecbf{w}_{D_{pca}}$ corresponding to largest eigenvalues. Form projection matrix $\mat{W} = [\vecbf{w}_1, \dots, \vecbf{w}_{D_{pca}}]$.
    \item Transformed data: $\mat{X}^{(4)} = \mat{X}^{(3)} \mat{W}$.
    \item Output: $\mat{X}^{(4)}$.
\end{itemize}

\end{document}